% !TeX program = xelatex
% SPDX-License-Identifier: MIT
% Copyright (c) 2026 Ziyu Peng and 刘毅涵
% Author: https://github.com/pzy54154631-hub?tab=repositories
\documentclass[UTF8,a4paper,fontset=fandol]{ctexart}
\usepackage[margin=2.5cm]{geometry}
\usepackage{amsmath,amssymb}
\usepackage[hidelinks]{hyperref}
\title{微积分学习卡片}
\author{Ziyu Peng \and 刘毅涵}
\date{}
\begin{document}
\maketitle
\section{从函数到导数}
设 \(f(x)=x^2+2x\)，其中 \(x\in\mathbb{R}\)。
\begin{align*}
  f'(x) &= 2x+2,\\
  f''(x) &= 2.
\end{align*}
在 \(x=-1\) 处，一阶导数为零；
又因二阶导数为正，此处为严格局部极小值点。

\section{积分与验证}
\begin{align*}
  \int f(x)\,\mathrm{d}x
    &= \frac{x^3}{3}+x^2+C,\\
  \int_0^1 f(x)\,\mathrm{d}x
    &= \left[\frac{x^3}{3}+x^2\right]_0^1
     = \frac{4}{3}.
\end{align*}
对原函数再求导，可以检查不定积分的结果。

\section{拓展到两个变量}
设 \(g(x,y)=x^2+xy+y^2\)，则
\[
\nabla g=
\begin{pmatrix}
  2x+y\\
  x+2y
\end{pmatrix}.
\]
本例统一把梯度写成列向量。
\end{document}
